\documentclass[12pt, letterpaper]{article}

\usepackage[utf8]{inputenc}
\usepackage[T1]{fontenc}
\usepackage[spanish]{babel}
\usepackage{geometry}
\usepackage{amsmath, amssymb, amsthm}
\usepackage{booktabs}
\usepackage{array}
\usepackage{microtype}
\usepackage{setspace}
\usepackage{parskip}
\usepackage{titlesec}
\usepackage{fancyhdr}
\usepackage{enumitem}
\usepackage{bm}
\usepackage{mathtools}
\usepackage{tcolorbox}
\usepackage{hyperref}

\tcbuselibrary{skins, breakable}

\geometry{top=2.8cm, bottom=3.0cm, left=3.2cm, right=3.2cm, headheight=14pt}
\setstretch{1.20}

\pagestyle{fancy}
\fancyhf{}
\fancyhead[C]{\small\textsc{Econometría --- Gujarati --- Ejercicio 13.13}}
\fancyfoot[C]{\small\thepage}
\renewcommand{\headrulewidth}{0.4pt}
\renewcommand{\footrulewidth}{0pt}

\titleformat{\section}
  {\large\bfseries\scshape}{\thesection.}{0.6em}{}
  [\vspace{-4pt}\rule{\linewidth}{0.4pt}]
\titleformat{\subsection}{\normalsize\bfseries}{\thesubsection.}{0.5em}{}
\titleformat{\subsubsection}{\normalsize\itshape}{}{0em}{}

\newtcolorbox{clave}{
  colback=white, colframe=black,
  boxrule=0.5pt, arc=3pt,
  left=8pt, right=8pt, top=6pt, bottom=6pt,
  breakable
}

\newcommand{\E}{\operatorname{E}}
\newcommand{\Cov}{\operatorname{Cov}}
\newcommand{\Var}{\operatorname{Var}}

\begin{document}

\begin{center}
  {\LARGE\bfseries Evaluación Crítica del Diagnóstico de Leamer}\\[6pt]
  {\LARGE\bfseries sobre la Brecha entre Teoría y Práctica Econométrica}\\[12pt]
  {\large\itshape ``Sacerdotes'' y ``pecadores'': metáfora, fundamento algebraico}\\
  {\large\itshape y propuesta de reforma metodológica}\\[14pt]
  \rule{0.55\linewidth}{0.6pt}\\[8pt]
  {\normalsize Ejercicio~13.13 --- \textit{Econometría}, Gujarati \& Porter, 5.a ed.}\\[4pt]
  {\small\textsc{Análisis conceptual, teórico y algebraico}}
\end{center}

\vspace{10pt}

%------------------------------------------------------------------------------
\section{El diagnóstico de Leamer: planteamiento y metáfora}
%------------------------------------------------------------------------------

Edward Leamer describe una división profesional que, a su juicio, neutraliza
el rigor de la econometría como disciplina científica. A un lado se sitúa el
\textbf{``sacerdocio célibe de teóricos estadísticos''}: quienes construyen
el edificio formal del Modelo Clásico de Regresión Lineal (MCRL) bajo
supuestos que difícilmente se cumplen con datos reales. Al otro, la
\textbf{``legión de incorregibles pecadores analistas de datos''}: quienes
violan sistemáticamente esos supuestos, lo saben, y se limitan a
\emph{confesar} sus infracciones en notas a pie de página sin modificar sus
conclusiones.

La metáfora religiosa no es ornamental. Leamer apunta a una disfunción
estructural: el sistema de incentivos de la profesión premia la publicación
de resultados estadísticamente significativos, no la honestidad metodológica.
El resultado es que la inferencia econométrica publicada refleja, en gran
medida, un proceso de ensayo y error no documentado, y que los valores $p$
reportados no tienen el significado probabilístico que la teoría les atribuye.
El objetivo de Leamer es reformar esta práctica, no destruir la econometría:
en sus propias palabras, se trata de
\emph{``quitarle el engaño a la econometría''}
(\textit{take the con out of econometrics}).

%------------------------------------------------------------------------------
\section{Fundamento teórico: por qué la brecha es estructural}
%------------------------------------------------------------------------------

\subsection{La naturaleza observacional de los datos económicos}

A diferencia de las ciencias experimentales, el economista no controla el
proceso generador de datos. Los problemas de \textbf{multicolinealidad},
\textbf{endogeneidad}, \textbf{errores de medición} y
\textbf{omisión de variables} no son anomalías evitables: son consecuencias
inevitables de trabajar con registros históricos donde ninguna variable
fue asignada aleatoriamente. Esta asimetría entre la situación ideal que
presupone la teoría y la realidad de los datos es la raíz conceptual de
la brecha denunciada por Leamer.

\subsection{Minería de datos y colapso de la inferencia clásica}

El ``pecado'' más extendido es la \textbf{minería de datos}
(\textit{data mining}): el analista estima múltiples especificaciones
del mismo modelo hasta encontrar una con coeficientes significativos
y signos económicamente ``razonables'', y luego reporta esa
especificación como si hubiera sido la única estimada. Esta práctica
colapsa la validez de las pruebas $t$ y $F$ clásicas, que están
construidas bajo el supuesto de que se estimó exactamente el modelo
pre-especificado, sin búsqueda previa.

\subsection{Significancia estadística \textit{versus} relevancia económica}

Un tercer pilar de la crítica es la confusión entre significancia
estadística y relevancia económica. Un coeficiente puede ser
estadísticamente distinto de cero gracias a un tamaño muestral
grande, pero su magnitud puede ser económicamente despreciable. La
obsesión profesional con los valores $p$ inferiores a $0{,}05$ desvía
la atención de la pregunta sustantiva: ¿es el efecto estimado lo
suficientemente grande como para importar en términos de política
económica o decisión empresarial?

%------------------------------------------------------------------------------
\section{Fundamento algebraico: el costo del ``pecado'' en la inferencia}
%------------------------------------------------------------------------------

\subsection{Ajuste de Lovell: inflación del error de Tipo I por minería de datos}

Sea $c$ el número total de variables candidatas examinadas y $k$ el número
de variables finalmente retenidas en el modelo, con $k < c$. Si el nivel
nominal de significancia es $\alpha$, la probabilidad real de cometer un
error de Tipo I (rechazar una hipótesis nula verdadera) ya no es $\alpha$
sino $\alpha^*$, dada por la \textbf{fórmula de Lovell}:
\begin{equation}
  \alpha^* = 1 - (1 - \alpha)^{c/k} \approx \frac{c}{k}\,\alpha.
  \label{eq:lovell}
\end{equation}

La aproximación de primer orden en~\eqref{eq:lovell} es válida cuando
$\alpha$ es pequeño. El resultado es contundente: si un analista examina
$c = 15$ variables para retener $k = 5$ a un nivel nominal de $\alpha = 0{,}05$,
el nivel real es:
\[
  \alpha^* \approx \frac{15}{5} \times 0{,}05 = 0{,}15.
\]
Es decir, la probabilidad de concluir erróneamente que una variable es
significativa es del $15\,\%$, tres veces mayor que el $5\,\%$ reportado.
La ``confesión'' (reportar el número de especificaciones exploradas) es
necesaria precisamente porque sin ella el lector no puede calcular $\alpha^*$
y atribuye a los resultados una fiabilidad que no poseen.

\subsection{Sesgo por omisión de variables: formalización del ``pecado'' de omisión}

Supóngase que el modelo verdadero es:
\begin{equation}
  Y_i = \beta_1 + \beta_2 X_{2i} + \beta_3 X_{3i} + u_i,
  \label{eq:verdadero}
\end{equation}
pero el analista, por falta de datos o descuido, estima:
\begin{equation}
  Y_i = \alpha_1 + \alpha_2 X_{2i} + v_i, \qquad v_i = \beta_3 X_{3i} + u_i.
  \label{eq:omitido}
\end{equation}
Como fue demostrado en el Ejercicio~13.12, los estimadores MCO del modelo
mal especificado satisfacen:
\begin{align}
  \E(\hat{\alpha}_2) &= \beta_2 + \beta_3\,b_{32}, \label{eq:sesgo_pend}\\
  \E(\hat{\alpha}_1) &= \beta_1 + \beta_3\,(\bar{X}_3 - b_{32}\,\bar{X}_2),
  \label{eq:sesgo_int}
\end{align}
donde $b_{32}$ es el coeficiente de la regresión auxiliar de $X_3$ sobre
$X_2$. Ambos estimadores son \textbf{sesgados e inconsistentes}, y el sesgo
persiste en muestras arbitrariamente grandes. Si el analista reporta
$\hat{\alpha}_2$ como estimación de $\beta_2$ sin advertir la omisión,
presenta como parámetro estructural un valor que confunde el efecto
directo de $X_2$ con la correlación espuria transmitida a través de la
variable excluida.

\subsection{Interacción entre los dos ``pecados''}

Los efectos de~\eqref{eq:lovell} y \eqref{eq:sesgo_pend}--\eqref{eq:sesgo_int}
se potencian mutuamente. El analista que minea datos tiende a retener
las variables cuyo sesgo $\beta_3 b_{32}$ empuja el coeficiente en la
dirección estadísticamente significativa, mientras descarta las variables
cuya inclusión haría evidente la magnitud del sesgo. El proceso de
selección sesgada de variables produce, por tanto, no solo un nivel de
significancia inflado, sino estimadores con sesgo de signo sistemático.

%------------------------------------------------------------------------------
\section{Evaluación crítica del diagnóstico de Leamer}
%------------------------------------------------------------------------------

\subsection{Aspectos en los que el diagnóstico es acertado}

\begin{enumerate}[leftmargin=1.8cm, label=\textbf{\arabic*.}]

  \item \textbf{Validez empírica de la metáfora.} La división que describe
    Leamer es verificable en la literatura: existe una producción abundante
    de teoría econométrica sofisticada y, simultáneamente, una práctica
    empírica que raramente aplica más de las pruebas diagnósticas básicas.
    La asimetría no es accidental; responde a incentivos institucionales
    concretos (revistas, contrataciones, citaciones).

  \item \textbf{La inflación del error Tipo I es un hecho matemático.}
    La fórmula de Lovell~\eqref{eq:lovell} no es una opinión: es un
    resultado probabilístico exacto. Que la profesión lo ignore
    sistemáticamente no lo hace menos válido.

  \item \textbf{El llamado a la transparencia es metodológicamente correcto.}
    Reportar el número de especificaciones exploradas, presentar análisis
    de sensibilidad ante cambios en la selección de variables, y distinguir
    entre significancia estadística y relevancia económica son prácticas
    que mejoran objetivamente la calidad de la inferencia.

\end{enumerate}

\subsection{Limitaciones y matices del diagnóstico}

\begin{enumerate}[leftmargin=1.8cm, label=\textbf{\arabic*.}]

  \item \textbf{El diagnóstico no ofrece una solución operativa completa.}
    Leamer propone el \textbf{análisis de sensibilidad extrema}
    (\textit{Extreme Bounds Analysis}, EBA) como alternativa: examinar
    cómo varía el coeficiente de interés al permutar el subconjunto de
    variables de control. Sin embargo, la EBA ha sido criticada por ser
    excesivamente conservadora (declara ``frágiles'' resultados que otras
    metodologías consideran robustos) y por carecer de fundamento
    probabilístico formal.

  \item \textbf{La metáfora oscurece la heterogeneidad interna de la
    práctica.} No todo análisis empírico es minería de datos irresponsable.
    La econometría estructural, los experimentos naturales, los diseños
    de regresión discontinua y los métodos de variables instrumentales
    representan avances sustanciales en la credibilidad empírica que
    Leamer no contemplaba en su diagnóstico original de 1983.

  \item \textbf{El estándar de ``robustez total'' puede ser
    inalcanzable.} En economía, como en cualquier ciencia social, la
    especificación del modelo es inevitablemente teoría-dependiente.
    Exigir que un resultado sea invariante ante \emph{cualquier}
    combinación de variables de control puede equivaler a exigir que
    los datos hablen sin teoría, lo cual es epistemológicamente incoherente.

\end{enumerate}

%------------------------------------------------------------------------------
\section{Propuesta de reforma metodológica}
%------------------------------------------------------------------------------

La evaluación de Leamer sugiere un cambio de paradigma cuyo contenido
positivo puede sistematizarse en cuatro principios:

\begin{enumerate}[leftmargin=1.8cm, label=\textbf{P\arabic*.}]

  \item \textbf{Reconocimiento de la subjetividad del modelado.} La
    especificación econométrica es simultáneamente una hipótesis científica
    y una decisión de diseño. Tratarla como un procedimiento puramente
    mecánico oscurece la responsabilidad intelectual del investigador.

  \item \textbf{Transparencia en el proceso de búsqueda.} El investigador
    debe reportar el número de especificaciones exploradas para permitir
    al lector aplicar la corrección de Lovell~\eqref{eq:lovell} y evaluar
    correctamente la fiabilidad de los resultados.

  \item \textbf{Análisis de sensibilidad sistemático.} En lugar de
    presentar un único modelo como ``el correcto'', se deben reportar
    los rangos de variación de los coeficientes clave ante cambios razonables
    en la especificación (inclusión o exclusión de controles, cambios
    en la forma funcional, submuestras temporales o geográficas).

  \item \textbf{Distinción explícita entre significancia estadística y
    relevancia económica.} Un resultado válido debe superar ambos umbrales:
    ser estadísticamente distinguible del ruido \emph{y} tener una
    magnitud sustantiva desde la perspectiva de la teoría económica o
    la política pública.

\end{enumerate}

\begin{clave}
  \textbf{Conclusión:} El diagnóstico de Leamer es esencialmente correcto
  en su descripción del problema y en su llamado a la honestidad intelectual.
  La fórmula de Lovell~\eqref{eq:lovell} y el sesgo por variable omitida
  [\eqref{eq:sesgo_pend}--\eqref{eq:sesgo_int}] demuestran algebraicamente
  que las ``confesiones'' son necesarias porque los resultados publicados
  sin ellas son estadística y econométricamente inválidos.
  La limitación del diagnóstico es que subestima los avances metodológicos
  posteriores (identificación causal, experimentos naturales, diseños
  quasi-experimentales) y que su propuesta operativa (EBA) no resolvió
  de manera satisfactoria el problema que identificó. El legado duradero
  de Leamer es haber convertido la \emph{fragilidad de la inferencia} en
  un objeto de estudio legítimo dentro de la econometría aplicada.
\end{clave}

\end{document}
